\documentclass[10pt,a4paper]{amsart}
\usepackage[margin=19mm]{geometry}
\usepackage{amsmath,amssymb,mathtools}
\usepackage[scr=rsfs,cal=euler]{mathalfa}
\usepackage{xfrac}
\usepackage[unicode,hidelinks,bookmarks=false]{hyperref}
\newcommand{\ie}{i.e.\ }
\newcommand{\cf}{cf.\ }
\newcommand{\resp}{respectively\ }
\newcommand{\Subsection}{\subsection}
\providecommand{\nobreakdash}{\nobreak\hbox{-}}
\setlength{\emergencystretch}{2em}
\multlinegap0pt
\usepackage{siunitx}
\usepackage{siunitx}
\usepackage{siunitx}
\usepackage{siunitx}
\usepackage{bm}
\usepackage{xcolor}
\usepackage[shortlabels]{enumitem}
\usepackage{tcolorbox}
\usepackage[normalem]{ulem}
\usepackage{amssymb}
\usepackage{siunitx}
\usepackage{algorithm}\usepackage{algpseudocode}
\newtheorem{theorem}{Theorem}
\DeclareMathOperator{\Cov}{Cov}
\DeclareMathOperator{\E}{\mathbb{E}}
\title{Gaussian Process Differential Ensembles for Joint Inference on Curves, Derivatives, and Integrals}
\author{Andreas Kryger Jensen\, and Adam Gorm Hoffmann}
\date{Source: arXiv:2606.23036v1 (CC BY 4.0)}
\begin{document}
\maketitle

\noindent\emph{Source and licence.} Gaussian Process Differential Ensembles for Joint Inference on Curves, Derivatives, and Integrals. Version arXiv:2606.23036v1 (CC BY 4.0), distributed by arXiv under the Creative Commons Attribution 4.0 International licence, http://creativecommons.org/licenses/by/4.0/, as recorded in the arXiv licence metadata for this exact version and shown on the abstract page https://arxiv.org/abs/2606.23036v1. Full licence notice: \texttt{licenses/arxiv-2606.23036-CC-BY-4.0.txt}.\par\smallskip

\noindent\textit{Editorial scope.} Counted unit: the article's theorem thm:appendix-schur (source0.tex lines 1848--1869). The article's complete original proof of it is delivered with it (source0.tex lines 1871--1906, one \texttt{proof} environment of 36 lines). No mathematics is written, repaired or re-derived: the statement and the proof are the article's own lines, in the article's own notation, and no retained line is substituted or re-flowed.  No further source-local context is needed: every cross-reference of every retained line resolves inside the two delivered spans.  Nothing was omitted from any retained span, so the package carries no omission record.  No retained line contains a citation, so the wrapper carries no bibliography and no bibliography entry.  No retained line contains a figure environment, an image-inclusion command or drawing code, so no image is delivered and the package declares no asset dependency.  Editorial formatting only, in addition: an AMS \texttt{amsart} preamble that restates the article's own declarations and macros used by the retained lines, namely the environments \texttt{theorem} on separate counters, together with the standard packages those lines need.  The retained environments print in this document as Theorem 1 in order of appearance, whereas the article prints the same items with the numbering of its own full document.  The complete native source of the article as distributed in the arXiv e-print for this exact version is bundled unchanged as \texttt{sources/203403/source0.tex}.
\begin{theorem}[Covariance function criterion for the dependent extension]\label{thm:appendix-schur}
Let \(\boldsymbol{\Sigma}_{\kappa}\in\mathbb R^{r\times r}\) be symmetric positive definite, and define
\[
C_{\theta\mid\kappa}(s,t):=C_\theta(s,t)-\boldsymbol\gamma(s)^\top\boldsymbol{\Sigma}_{\kappa}^{-1}\boldsymbol{\gamma}(t).
\]
Then the following are equivalent:
\begin{enumerate}
\item There exists a jointly Gaussian pair \((f_0,\boldsymbol{\kappa})\) such that
\[
\E[f_0(t)]=\mu(t),\quad t\in[t_0,t_1],
\qquad
\E[\boldsymbol{\kappa}]=\boldsymbol{\mu}_\kappa,
\]
and
\begin{align*}
\Cov(f_0(s),f_0(t))&=C_\theta(s,t), && s,t\in[t_0,t_1],\\
\Cov(f_0(s),\boldsymbol{\kappa})&=\boldsymbol{\gamma}(s), && s\in[t_0,t_1],\\
\Cov(\boldsymbol{\kappa})&=\boldsymbol{\Sigma}_{\kappa}.
\end{align*}
\item The covariance function \(C_{\theta\mid\kappa}\) is positive semidefinite on \([t_0,t_1]\), that is, for every \(n\in\mathbb{N}\) and every \(u_1,\ldots,u_n\in[t_0,t_1]\), the matrix \(\bigl(C_{\theta\mid\kappa}(u_i,u_j)\bigr)_{i,j=1}^n\) is positive semidefinite.
\end{enumerate}
\end{theorem}

\begin{proof}
Assume first that $(f_0,\boldsymbol{\kappa})$ defines a valid jointly Gaussian model. Then, for every $n\in\mathbb{N}$ and every $u_1,\ldots,u_n\in[t_0,t_1]$, the matrix
\begin{align*}
\boldsymbol{\Omega}_n=
\begin{pmatrix}
\mathbf{C}_n & \mathbf{G}_n\\
\mathbf{G}_n^\top & \boldsymbol{\Sigma}_{\kappa}
\end{pmatrix}
\end{align*}
is positive semidefinite. Since $\boldsymbol{\Sigma}_{\kappa}$ is positive definite, the standard Schur complement criterion for positive semidefinite block matrices yields
\begin{align*}
\boldsymbol{\Omega}_n\succeq 0
\quad\Longleftrightarrow\quad
\mathbf{C}_n-\mathbf{G}_n\boldsymbol{\Sigma}_{\kappa}^{-1}\mathbf{G}_n^\top\succeq 0.
\end{align*}
The finite-dimensional admissibility condition is therefore Schur-complement positivity.
The $(i,j)$ entry of the Schur complement is
\begin{align*}
C_\theta(u_i,u_j)-\boldsymbol{\gamma}(u_i)^\top\boldsymbol{\Sigma}_{\kappa}^{-1}\boldsymbol{\gamma}(u_j)
=
C_{\theta\mid\kappa}(u_i,u_j).
\end{align*}
Hence the matrix $\bigl(C_{\theta\mid\kappa}(u_i,u_j)\bigr)_{i,j=1}^n$ is positive semidefinite for every finite choice of points, so $C_{\theta\mid\kappa}$ is positive semidefinite on $[t_0,t_1]$.

Conversely, assume that $C_{\theta\mid\kappa}$ is positive semidefinite on $[t_0,t_1]$. Then for every $n\in\mathbb{N}$ and every $u_1,\ldots,u_n\in[t_0,t_1]$, the matrix
\begin{align*}
\mathbf{C}_n-\mathbf{G}_n\boldsymbol{\Sigma}_{\kappa}^{-1}\mathbf{G}_n^\top
=
\bigl(C_{\theta\mid\kappa}(u_i,u_j)\bigr)_{i,j=1}^n
\end{align*}
is positive semidefinite. Since $\boldsymbol{\Sigma}_{\kappa}$ is positive definite, the Schur complement criterion implies that $\boldsymbol{\Omega}_n\succeq 0$ for every finite set of evaluation points. Assign to
\[
\bigl(f_0(u_1),\ldots,f_0(u_n),\boldsymbol{\kappa}^\top\bigr)^\top
\]
the Gaussian law with the prescribed mean vector and covariance matrix $\boldsymbol{\Omega}_n$. These finite-dimensional Gaussian laws are consistent under deletion and permutation of evaluation points because the prescribed means and covariances restrict to the corresponding sub-vectors and principal submatrices. Kolmogorov's extension theorem then supplies a process on the product index set \([t_0,t_1]\) together with the \(r\) distinguished coordinates of \(\boldsymbol{\kappa}\). Since every finite subvector was assigned a Gaussian law, the resulting pair \((f_0,\boldsymbol{\kappa})\) is jointly Gaussian with the displayed mean and covariance blocks.
\end{proof}

\end{document}
